\documentclass[12pt,letterpaper]{article}
\usepackage{graphicx,textcomp}
\usepackage{natbib}
\usepackage{setspace}
\usepackage{fullpage}
\usepackage{color}
\usepackage[reqno]{amsmath}
\usepackage{amsthm}
\usepackage{fancyvrb}
\usepackage{amssymb,enumerate}
\usepackage[all]{xy}
\usepackage{endnotes}
\usepackage{lscape}
\newtheorem{com}{Comment}
\usepackage{float}
\usepackage{hyperref}
\newtheorem{lem} {Lemma}
\newtheorem{prop}{Proposition}
\newtheorem{thm}{Theorem}
\newtheorem{defn}{Definition}
\newtheorem{cor}{Corollary}
\newtheorem{obs}{Observation}
\usepackage[compact]{titlesec}
\usepackage{dcolumn}
\usepackage{tikz}
\usetikzlibrary{arrows}
\usepackage{multirow}
\usepackage{xcolor}
\newcolumntype{.}{D{.}{.}{-1}}
\newcolumntype{d}[1]{D{.}{.}{#1}}
\definecolor{light-gray}{gray}{0.65}
\usepackage{url}
\usepackage{listings}
\usepackage{color}

\definecolor{codegreen}{rgb}{0,0.6,0}
\definecolor{codegray}{rgb}{0.5,0.5,0.5}
\definecolor{codepurple}{rgb}{0.58,0,0.82}
\definecolor{backcolour}{rgb}{0.95,0.95,0.92}

\lstdefinestyle{mystyle}{
	backgroundcolor=\color{backcolour},   
	commentstyle=\color{codegreen},
	keywordstyle=\color{magenta},
	numberstyle=\tiny\color{codegray},
	stringstyle=\color{codepurple},
	basicstyle=\footnotesize,
	breakatwhitespace=false,         
	breaklines=true,                 
	captionpos=b,                    
	keepspaces=true,                 
	numbers=left,                    
	numbersep=5pt,                  
	showspaces=false,                
	showstringspaces=false,
	showtabs=false,                  
	tabsize=2
}
\lstset{style=mystyle}
\newcommand{\Sref}[1]{Section~\ref{#1}}
\newtheorem{hyp}{Hypothesis}


\title{Problem Set 4}
\date{Sein Kim}
\author{Applied Stats/Quant Methods 1}


\begin{document}
	\maketitle


	\vspace{.5cm}
\section*{Question 1: Economics}
\vspace{.25cm}
\noindent 	
In this question, use the \texttt{prestige} dataset in the \texttt{car} library. First, run the following commands:

\begin{verbatim}
install.packages(car)
library(car)
data(Prestige)
help(Prestige)
\end{verbatim} 


\noindent We would like to study whether individuals with higher levels of income have more prestigious jobs. Moreover, we would like to study whether professionals have more prestigious jobs than blue and white collar workers.

\begin{enumerate}
	
	\item [(a)]
	Create a new variable \texttt{professional} by recoding the variable \texttt{type} so that professionals are coded as $1$, and blue and white collar workers are coded as $0$ (Hint: \texttt{ifelse}).
	
	
	\noindent\textbf{R Code:}
	\lstinputlisting[language=R, firstline=40, lastline=43]{PS04_SK.R}
	
	\noindent\textbf{Answer:}\\
	I create a dummy variable \texttt{professional} such that
	\[
	\texttt{professional}_i = 
	\begin{cases}
		1, & \text{if } \texttt{type}_i = \text{``prof'' (professional)}\\
		0, & \text{if } \texttt{type}_i = \text{``bc'' or ``wc''} \text{ (blue/white collar)}.
	\end{cases}
	\]
	Thus, \texttt{professional} indicates whether an observation belongs to a professional occupation.
	
	\vspace{0.4cm}
	
	\item [(b)]
	Run a linear model with \texttt{prestige} as an outcome and \texttt{income}, \texttt{professional}, and the interaction of the two as predictors (Note: this is a continuous $\times$ dummy interaction.)
	
	  \noindent\textbf{R Code:}
	\lstinputlisting[language=R, firstline=46, lastline=50]{PS04_SK.R}
	
	\noindent\textbf{Result:}\\
	I estimate the following linear regression model with an interaction between a continuous predictor and a dummy:
	\[
	\widehat{\text{prestige}} 
	= \hat{\beta}_0 
	+ \hat{\beta}_1\,\texttt{income}
	+ \hat{\beta}_2\,\texttt{professional}
	+ \hat{\beta}_3\,(\texttt{income} \times \texttt{professional}).
	\]
	The regression output reports coefficient estimates, standard errors, $t$-statistics, and $p$-values for each term, as well as the model $R^2$ and sample size.
	
	\vspace{0.4cm}
	
	\item [(c)]
	Write the prediction equation based on the result.
	
	
	\noindent\textbf{Prediction Equation:}\\
	Using the notation above, the fitted model can be written as
	\[
	\widehat{\text{prestige}} 
	= \hat{\beta}_0 
	+ \hat{\beta}_1\,\texttt{income}
	+ \hat{\beta}_2\,\texttt{professional}
	+ \hat{\beta}_3\,(\texttt{income} \times \texttt{professional}).
	\]
	For non-professionals (\texttt{professional} $=0$),
	\[
	\widehat{\text{prestige}}_{\text{non-prof}} 
	= \hat{\beta}_0 + \hat{\beta}_1\,\texttt{income},
	\]
	and for professionals (\texttt{professional} $=1$),
	\[
	\widehat{\text{prestige}}_{\text{prof}} 
	= (\hat{\beta}_0 + \hat{\beta}_2) + (\hat{\beta}_1 + \hat{\beta}_3)\,\texttt{income}.
	\]
	
	\vspace{0.4cm}
	
	\item [(d)]
	Interpret the coefficient for \texttt{income}.
	
	\noindent
	  In this interaction model, $\hat{\beta}_1$ is the marginal effect of income on the prestige score for \emph{non-professional} occupations (i.e., when \texttt{professional} $= 0$). 
	Holding \texttt{professional} fixed at $0$, a one-unit increase in \texttt{income} is associated with an expected change of $\hat{\beta}_1$ units in the prestige score. 
	In other words, among blue and white collar workers, higher income is associated with higher occupational prestige at a rate given by $\hat{\beta}_1$.
	
	\vspace{0.4cm}
	\newpage	
	\item [(e)]
	Interpret the coefficient for \texttt{professional}.
	
	  The coefficient $\hat{\beta}_2$ represents the difference in the expected prestige score between professional and non-professional occupations \emph{when income is equal to zero}. Formally,
	\[
	\widehat{\text{prestige}}_{\text{prof}} - \widehat{\text{prestige}}_{\text{non-prof}}
	= \hat{\beta}_2 \quad \text{when } \texttt{income} = 0.
	\]
	Since an income of zero is not substantively realistic in this context, $\hat{\beta}_2$ mainly captures a vertical shift in the intercept and is not directly interpreted as a meaningful average difference in prestige without also considering the interaction term.
	
	\vspace{0.4cm}
	
	\noindent 	
	\item [(f)]
	What is the effect of a \$1,000 increase in income on prestige score for professional occupations? In other words, we are interested in the marginal effect of income when the variable \texttt{professional} takes the value of $1$. Calculate the change in $\hat{y}$ associated with a \$1,000 increase in income based on your answer for (c).
	
	\noindent
	For professionals (\texttt{professional} $=1$), the marginal effect of income on prestige is
	\[
	\frac{\partial \widehat{\text{prestige}}}{\partial\,\texttt{income}}
	= \hat{\beta}_1 + \hat{\beta}_3.
	\]
	Therefore, a \$1,000 increase in income changes the predicted prestige score by
	\[
	\Delta \widehat{\text{prestige}}
	= 1{,}000 \times (\hat{\beta}_1 + \hat{\beta}_3).
	\]
	Plugging in the estimated values of $\hat{\beta}_1$ and $\hat{\beta}_3$ from the regression output gives the numerical change in the prestige score associated with a \$1,000 income increase for professional occupations.
	\vspace{0.4cm}
	
	\item [(g)]
	What is the effect of changing one's occupations from non-professional to professional when her income is \$6,000? We are interested in the marginal effect of professional jobs when the variable \texttt{income} takes the value of $6,000$. Calculate the change in $\hat{y}$ based on your answer for (c).
	
	\noindent   At a fixed income level of $\texttt{income} = 6{,}000$, the predicted prestige score for non-professionals is
	\[
	\widehat{\text{prestige}}_{\text{non-prof}} 
	= \hat{\beta}_0 + \hat{\beta}_1 \times 6{,}000,
	\]
	and for professionals,
	\[
	\widehat{\text{prestige}}_{\text{prof}} 
	= (\hat{\beta}_0 + \hat{\beta}_2) + (\hat{\beta}_1 + \hat{\beta}_3)\times 6{,}000.
	\]
	The marginal effect of switching from non-professional to professional at income \$6,000 is therefore
	\[
	\Delta \widehat{\text{prestige}}
	= \widehat{\text{prestige}}_{\text{prof}} - \widehat{\text{prestige}}_{\text{non-prof}}
	= \hat{\beta}_2 + 6{,}000\,\hat{\beta}_3.
	\]
	Using the estimated coefficients $\hat{\beta}_2$ and $\hat{\beta}_3$ from the regression, this expression yields the predicted change in prestige when an individual with income \$6,000 moves from a non-professional to a professional occupation.
	
\end{enumerate}

\newpage

\section*{Question 2: Political Science}
\vspace{.25cm}
\noindent 	Researchers are interested in learning the effect of all of those yard signs on voting preferences.\footnote{Donald P. Green, Jonathan	S. Krasno, Alexander Coppock, Benjamin D. Farrer,	Brandon Lenoir, Joshua N. Zingher. 2016. ``The effects of lawn signs on vote outcomes: Results from four randomized field experiments.'' Electoral Studies 41: 143-150. } Working with a campaign in Fairfax County, Virginia, 131 precincts were randomly divided into a treatment and control group. In 30 precincts, signs were posted around the precinct that read, ``For Sale: Terry McAuliffe. Don't Sellout Virgina on November 5.'' \\

Below is the result of a regression with two variables and a constant.  The dependent variable is the proportion of the vote that went to McAuliff's opponent Ken Cuccinelli. The first variable indicates whether a precinct was randomly assigned to have the sign against McAuliffe posted. The second variable indicates
a precinct that was adjacent to a precinct in the treatment group (since people in those precincts might be exposed to the signs).  \\

\vspace{.5cm}
\begin{table}[!htbp]
	\centering 
	\textbf{Impact of lawn signs on vote share}\\
	\begin{tabular}{@{\extracolsep{5pt}}lccc} 
		\\[-1.8ex] 
		\hline \\[-1.8ex]
		Precinct assigned lawn signs  (n=30)  & 0.042\\
		& (0.016) \\
		Precinct adjacent to lawn signs (n=76) & 0.042 \\
		&  (0.013) \\
		Constant  & 0.302\\
		& (0.011)
		\\
		\hline \\
	\end{tabular}\\
	\footnotesize{\textit{Notes:} $R^2$=0.094, N=131}
\end{table}

\vspace{.5cm}
\begin{enumerate}
  \item[(a)] Use the results from a linear regression to determine whether having these yard signs in a precinct affects vote share (e.g., conduct a hypothesis test with $\alpha = .05$).\\[0.2cm]

The coefficient on ``Precinct assigned lawn signs'' is $\hat{\beta}_1 = 0.042$ with a standard error of $0.016$. The corresponding $t$-statistic is
\[
t = \frac{0.042}{0.016} = 2.625.
\]
Using a two-sided $t$-test with $\alpha = 0.05$, the critical value is approximately $|t| \approx 1.98$ for large degrees of freedom. Since $|2.625| > 1.98$, we reject the null hypothesis $H_0 : \beta_1 = 0$ at the 5\% level. 

Substantively, assigning lawn signs to a precinct increases the opponent’s vote share by about 4.2 percentage points on average, and this effect is statistically significant at $\alpha = 0.05$.

\vspace{0.4cm}

	\item [(b)]  Use the results to determine whether being
	next to precincts with these yard signs affects vote
	share (e.g., conduct a hypothesis test with $\alpha = .05$). \\[0.2cm]
	
	\noindent   The coefficient on ``Precinct adjacent to lawn signs'' is also $\hat{\beta}_2 = 0.042$ with a standard error of $0.013$. The $t$-statistic is
	\[
	t = \frac{0.042}{0.013} \approx 3.23.
	\]
	Again using a two-sided test with $\alpha = 0.05$, this $t$-value exceeds the critical value in absolute terms, so we reject $H_0 : \beta_2 = 0$. 
	
	Substantively, precincts adjacent to those with lawn signs also see an increase of about 4.2 percentage points in the opponent’s vote share, and this spillover effect is statistically significant at the 5\% level.
	
	\vspace{0.4cm}
	
	\item [(c)] Interpret the coefficient for the constant term substantively. \\[0.2cm]
	 The constant term is estimated as $\hat{\beta}_0 = 0.302$. This represents the expected proportion of the vote going to Cuccinelli in precincts that are \emph{neither} assigned lawn signs nor adjacent to precincts with lawn signs. In other words, in control precincts with no direct or neighboring exposure to the signs, the model predicts that Cuccinelli receives about 30.2\% of the vote.
	
	\vspace{0.4cm}

	\item [(d)] Evaluate the model fit for this regression.  What does this	tell us about the importance of yard signs versus other factors that are not modeled? \\[0.2cm]
	\noindent
	The model’s $R^2$ is reported as $0.094$, which means that the two indicators for lawn signs (treatment and adjacency), together with the constant, explain about 9.4\% of the variation in Cuccinelli’s vote share across precincts. While the coefficients on the lawn sign variables are statistically significant, the overall fit is relatively low, indicating that most of the variation in vote share is left unexplained by this simple model.
	
	Substantively, this suggests that yard signs may have a real, detectable effect on voting outcomes, but many other factors—such as demographics, partisanship, campaign activities, and local political context—are likely more important drivers of precinct-level vote share than lawn signs alone.
	
\end{enumerate}  


\end{document}
